% 【本文件写什么】一级组件 wateros-runtime（§ 级聚合）
%   - 组件职责与在内核中的位置：控制台、分级日志、panic、堆分配器与串口再导出。
%   - 聚合 crate os/components/wateros-runtime/src/lib.rs 的 pub mod 树与对外导出面
%   - 子模块划分、与相邻组件的依赖边界（谁依赖谁、走哪条门面）
%   - 根 feature / 本组件 Cargo.toml feature 如何选用各 impl
%   - 1～2 段综述后，由下方 \input 展开子模块；不在此逐条列举 api trait
% 【事实来源】
%   - os/components/wateros-runtime/src/lib.rs、Cargo.toml
%   - docs/exports/public-api/wateros-runtime.md
%   - docs/exports/features/wateros-runtime.md
%   - os/feature-tree.txt 中 wateros-runtime 相关段
% 【不写什么】
%   - api-v0 契约细节 → 子目录 api/api-v0.tex
%   - 各 impl 算法/平台细节 → 子目录 impl/*.tex

\section{运行时组件}

运行时层\code{wateros-runtime}聚合panic处理、控制台输出、开发日志与内核堆分配等早期运行能力。根crate只做子模块再导出，不在此定义新的初始化语义；各子能力的行为由对应子crate的feature决定，根内核按引导顺序分别挂接panic入口、初始化分配器、注册日志器并打开控制台。

QEMU主线默认走平台控制台路径，把打印宏接到平台早期串口；占位feature下任何实际写控制台会触发未实现panic。开发日志经着色输出上屏展示，不写入内核日志环；需要同时上屏与持久化时，调用方须分别走日志宏与klog路径。日志级别由trace到error的一组feature择一，多开时取最安静档。

可选串口feature再导出QEMU\code{virt} UART驱动，供伪shell等路径使用，与控制台路径分工明确、互不替代。panic路径在着色打印后请求平台关机；控制台契约当前仅覆盖写侧，无输入API。内核堆默认使用链表分配器，也可选用TLSF后端；分配路径带中断屏蔽与递归检测，接近高水位时会发出警告。

\begin{rustcode}
// runtime/lib.rs: 子能力以 runtime::console 等形式再导出
pub mod panic {
    pub use panic::panic_handler; // 供根 crate 挂接 #[panic_handler]
}
pub mod console {
    pub use console::*;
}
pub mod logging {
    pub use logging::*;
}
pub mod heap_allocator {
    pub use heap_allocator::*;
}
#[cfg(feature = "serial-uart-virt")]
pub mod serial {
    pub use ::runtime_serial::*;
}
\end{rustcode}

\begin{rustcode}
// runtime-panic: 着色打印后请求平台关机
pub fn panic_handler(_panic_info: &core::panic::PanicInfo) -> ! {
    use console::{println, AnsiColor};
    if let Some(location) = _panic_info.location() {
        println!("... Panicked at {}:{}.  {}",
                 location.file(), location.line(), _panic_info.message());
    }
    use platform::reset::{shutdown, PlatformResetReason};
    while let Err(_e) = shutdown(PlatformResetReason::SystemFailure) {}
    loop {}
}
\end{rustcode}
